\documentclass[10.5pt,a4paper]{article}
% shared preamble for the English and Bulgarian papers (XeTeX; built with tectonic)
\usepackage[a4paper,left=26mm,right=26mm,top=26mm,bottom=28mm]{geometry}
\usepackage{fontspec}
\setmainfont{Palatino}
\setmonofont{Menlo}[Scale=MatchLowercase]
\usepackage{polyglossia}
\newfontfamily\cyrillicfont{Palatino}
\newfontfamily\cyrillicfonttt{Menlo}[Scale=MatchLowercase]
\usepackage{microtype}
\usepackage{booktabs,array,tabularx}
\usepackage{enumitem}
\usepackage{titlesec}
\usepackage{fancyhdr}
\usepackage{natbib}
\usepackage[hidelinks]{hyperref}
\usepackage{xurl}
\usepackage[font=small,labelfont=bf,labelsep=period,justification=raggedright,singlelinecheck=false,skip=4pt]{caption}

\usepackage{pgfplots}
\usepgfplotslibrary{groupplots}
\pgfplotsset{compat=1.17}
\definecolor{refcolor}{HTML}{0072B2}
\definecolor{varcolor}{HTML}{D55E00}
\definecolor{repcolor}{HTML}{009E73}

\setlength{\parindent}{0pt}
\setlength{\emergencystretch}{2em}
\setlength{\parskip}{5pt plus 1pt}
\linespread{1.07}
\bibpunct{(}{)}{;}{a}{}{,}
\setlength{\bibsep}{1.5pt}
\renewcommand{\bibfont}{\small}
\setlist{nosep,leftmargin=1.4em,itemsep=2pt}
\titleformat{\section}{\large\bfseries}{\thesection}{0.6em}{}
\titleformat{\subsection}{\normalsize\bfseries}{\thesubsection}{0.6em}{}
\titlespacing*{\section}{0pt}{14pt}{5pt}
\titlespacing*{\subsection}{0pt}{9pt}{3pt}
\pagestyle{plain}
\renewcommand{\arraystretch}{1.15}
\newcolumntype{L}[1]{>{\raggedright\arraybackslash}p{#1}}
\newcommand{\code}[1]{\texttt{\small #1}}

\newcommand{\titleblock}[3]{%
  \begin{center}
    {\Large\bfseries #1\par}\vspace{8pt}
    {\large #2\par}\vspace{2pt}
    {\small #3\par}
  \end{center}\vspace{4pt}}
\newenvironment{abstractbox}[1]{\begin{center}\bfseries\small #1\end{center}\vspace{-6pt}\begin{quote}\small}{\end{quote}}

% worked-example figure. arguments: time label, reference, edit inside the network, passive reporter, occupancy title.
% data: figure-patterning.csv, written from the Studio simulation of GAACACCCA (GLI2), signal 0.2 for 20 model units.
\newcommand{\patterningplot}[5]{%
\begin{tikzpicture}
\begin{groupplot}[group style={group size=3 by 2,horizontal sep=15mm,vertical sep=17mm},
  width=44mm,height=30mm,xmin=0,xmax=20,xtick={0,10,20},ymin=0,scaled y ticks=false,
  yticklabel style={/pgf/number format/fixed,/pgf/number format/precision=3},
  tick label style={font=\scriptsize},label style={font=\scriptsize},title style={font=\footnotesize\bfseries},
  every axis plot/.append style={thick}]
\nextgroupplot[title={Pax6}]
\addplot[refcolor] table[x=t,y=refPax6,col sep=comma]{figure-patterning.csv};
\addplot[varcolor,dashed] table[x=t,y=varPax6,col sep=comma]{figure-patterning.csv};
\nextgroupplot[title={Olig2}]
\addplot[refcolor] table[x=t,y=refOlig2,col sep=comma]{figure-patterning.csv};
\addplot[varcolor,dashed] table[x=t,y=varOlig2,col sep=comma]{figure-patterning.csv};
\nextgroupplot[title={Nkx2.2}]
\addplot[refcolor] table[x=t,y=refNkxTwoTwo,col sep=comma]{figure-patterning.csv};
\addplot[varcolor,dashed] table[x=t,y=varNkxTwoTwo,col sep=comma]{figure-patterning.csv};
\nextgroupplot[title={Irx3},xlabel={#1}]
\addplot[refcolor] table[x=t,y=refIrx3,col sep=comma]{figure-patterning.csv};
\addplot[varcolor,dashed] table[x=t,y=varIrx3,col sep=comma]{figure-patterning.csv};
\nextgroupplot[title={#5},xlabel={#1},ymax=1]
\addplot[refcolor] table[x=t,y=occRef,col sep=comma]{figure-patterning.csv};
\addplot[repcolor,densely dotted] table[x=t,y=occRep,col sep=comma]{figure-patterning.csv};
\addplot[varcolor,dashed] table[x=t,y=occVar,col sep=comma]{figure-patterning.csv};
\end{groupplot}
\node[anchor=north west,align=left,font=\scriptsize,inner sep=0pt] at ([xshift=12mm,yshift=-1mm]group c2r2.north east) {%
\tikz[baseline=-0.5ex]\draw[thick,refcolor](0,0)--(7mm,0);~#2\\[4pt]
\tikz[baseline=-0.5ex]\draw[thick,varcolor,dashed](0,0)--(7mm,0);~#3\\[4pt]
\tikz[baseline=-0.5ex]\draw[thick,repcolor,densely dotted](0,0)--(7mm,0);~#4};
\end{tikzpicture}}

\usetikzlibrary{arrows.meta}

% predictor pipeline. arguments: edit pair, window, 5-mer block, motif block, linear score, panel.
\newcommand{\pipelinefig}[6]{%
\begin{tikzpicture}[box/.style={draw,rounded corners=1.5pt,align=center,font=\scriptsize,inner sep=3pt,text width=24mm,minimum height=13mm},
  arr/.style={-{Stealth[length=1.8mm]},thick}]
\node[box] (a) at (0,0) {#1};
\node[box] (b) at (30mm,0) {#2};
\node[box] (c) at (60mm,10mm) {#3};
\node[box] (d) at (60mm,-10mm) {#4};
\node[box] (e) at (90mm,0) {#5};
\node[box] (f) at (120mm,0) {#6};
\draw[arr] (a)--(b);
\draw[arr] (b.east)--(c.west);
\draw[arr] (b.east)--(d.west);
\draw[arr] (c.east)--(e.west);
\draw[arr] (d.east)--(e.west);
\draw[arr] (e)--(f);
\end{tikzpicture}}

% internal signed utility by panel size (build/neural-predictor-reviewed-20261005/siraj-huber-joint-20261019/RESULTS.md).
% arguments: k5, joint-linear, Huber-linear, x label, y label.
\newcommand{\utilityplot}[5]{%
\begin{tikzpicture}
\begin{axis}[ybar,bar width=7pt,width=88mm,height=42mm,symbolic x coords={10,50,200},xtick=data,enlarge x limits=0.25,
  ymin=0,ymax=2.5,ylabel={#5},xlabel={#4},tick label style={font=\scriptsize},label style={font=\scriptsize},
  legend style={font=\scriptsize,at={(0.97,0.97)},anchor=north east,draw=none,fill=none},legend cell align=left,area legend]
\addplot[fill=black!35,draw=none] coordinates {(10,0.857) (50,0.607) (200,0.509)};
\addplot[fill=refcolor,draw=none] coordinates {(10,1.601) (50,0.954) (200,0.626)};
\addplot[fill=varcolor,draw=none] coordinates {(10,1.483) (50,0.986) (200,0.634)};
\legend{#1,#2,#3}
\end{axis}
\end{tikzpicture}}

% Scan of the 200-base example, from figure-scan.csv (docs/paper/figure-data.py). arguments: x label, colour-bar label, example-edit label.
\newcommand{\scanplot}[3]{%
\begin{tikzpicture}
\begin{axis}[width=124mm,height=24mm,scale only axis,axis on top,clip=false,y dir=reverse,xmin=0.5,xmax=200.5,ymin=-0.5,ymax=3.5,
  xtick={1,50,100,150,200},ytick={0,1,2,3},yticklabels={A,C,G,T},xlabel={#1},ylabel={},
  tick label style={font=\scriptsize},label style={font=\scriptsize},
  colormap={eff}{color(0cm)=(refcolor); color(1cm)=(white); color(2cm)=(varcolor)},
  point meta min=-0.6,point meta max=0.6,unbounded coords=jump,
  colorbar,colorbar style={width=2.5mm,ytick={-0.6,0,0.6},yticklabel style={font=\scriptsize,/pgf/number format/fixed,/pgf/number format/precision=1},
    ylabel={#2},ylabel style={font=\scriptsize}}]
\addplot[matrix plot*,mesh/cols=200,point meta=explicit] table[x=pos,y=altidx,meta=effect,col sep=comma]{figure-scan.csv};
\draw[-{Stealth[length=1.6mm]},thick] (axis cs:100,-1.9) -- (axis cs:100,-0.65);
\node[font=\scriptsize,anchor=south] at (axis cs:100,-1.9) {#3};
\end{axis}
\end{tikzpicture}}

% per-study correlation of the primary five studies, from figure-studies.csv. arguments: x label, Huber-linear, S2, k5.
\newcommand{\correlationplot}[4]{%
\begin{tikzpicture}
\begin{axis}[width=58mm,height=44mm,scale only axis,y dir=reverse,ymin=-0.6,ymax=4.6,ytick={0,1,2,3,4},yticklabels={Lee,McAfee,Matoba,Guo,Weiss},
  xmin=-0.05,xmax=0.15,xtick={-0.05,0,0.05,0.10,0.15},scaled x ticks=false,xticklabel style={/pgf/number format/fixed,/pgf/number format/precision=2},xlabel={#1},tick label style={font=\scriptsize},label style={font=\scriptsize},
  xmajorgrids,major grid style={black!12},extra x ticks={0},extra x tick style={grid=major,major grid style={black!55}},extra x tick labels={},
  legend style={font=\scriptsize,at={(0.5,1.04)},anchor=south,draw=none,legend columns=-1,/tikz/every even column/.append style={column sep=3mm}},legend cell align=left]
\addplot[only marks,mark=*,mark size=2.2pt,varcolor,restrict expr to domain={\thisrow{primary}}{1:1}] table[x=huber,y expr=\coordindex,col sep=comma]{figure-studies.csv};
\addplot[only marks,mark=square*,mark size=2pt,refcolor,restrict expr to domain={\thisrow{primary}}{1:1}] table[x=s2,y expr=\coordindex,col sep=comma]{figure-studies.csv};
\addplot[only marks,mark=triangle*,mark size=2.4pt,black!45,restrict expr to domain={\thisrow{primary}}{1:1}] table[x=k5,y expr=\coordindex,col sep=comma]{figure-studies.csv};
\legend{#2,#3,#4}
\end{axis}
\end{tikzpicture}}

% Huber-linear error relative to predicting no effect: internal tuning partition (0.28766713 against 0.30407514) and the five studies
% (figure-studies.csv; the mean bar is 0.10437080 against 0.09529676). arguments: x label, internal, five-study mean.
\newcommand{\excessplot}[3]{%
\begin{tikzpicture}
\begin{axis}[xbar,bar width=6pt,width=52mm,height=44mm,scale only axis,y dir=reverse,ymin=-0.6,ymax=6.6,ytick={0,1,2,3,4,5,6},
  yticklabels={#2,Lee,McAfee,Matoba,Guo,Weiss,#3},xmin=-10,xmax=60,xtick={0,20,40,60},xlabel={#1},
  tick label style={font=\scriptsize},label style={font=\scriptsize},xmajorgrids,major grid style={black!12},
  extra x ticks={0},extra x tick style={grid=major,major grid style={black!55}},extra x tick labels={}]
\addplot[fill=refcolor,draw=none] coordinates {(-5.40,0)};
\addplot[fill=varcolor,draw=none,restrict expr to domain={\thisrow{primary}}{1:1}] table[x=excess,y expr=\coordindex+1,col sep=comma]{figure-studies.csv};
\addplot[fill=black!55,draw=none] coordinates {(9.52,6)};
\end{axis}
\end{tikzpicture}}

\setmainlanguage{bulgarian}
\setotherlanguage{english}
\hypersetup{pdftitle={Epigenesis: промяна на ДНК и симулация на последствията за невронното развитие},pdfauthor={Деян Тодоров}}

\begin{document}

\titleblock{Epigenesis: промяна на ДНК и симулация\\на последствията за невронното развитие}{Деян Тодоров}{Технически доклад, чернова 2 · 20 октомври 2026 г. · English version: \href{epigenesis-technical-paper.pdf}{epigenesis-technical-paper.pdf}\\Код: \url{https://github.com/dttdrv/epigenesis} (Apache-2.0)}

\begin{abstractbox}{Резюме}
Epigenesis е софтуер, с който се променя участък от ДНК и се наблюдава какво следва в симулирана развиваща се система. Той чете ДНК източник при изрично тълкуване на значението на ДНК, компилира го в адресируем по съдържание модул за развитие (Development Module) и повтаря всяка компилация с отделен валидатор. Докладът описва изпълнимите компоненти и разглежда един пример докрай. Публикуваната замяна на GACCACCCA с GAACACCCA отслабва сайт за свързване на GLI в енхансер на \emph{Nkx2.2}. Измереното съотношение на свързване, 0,042, влиза в модел на оформяне на невралната тръба с четири гена. При постоянен сигнал 0,2 в продължение на 20 единици на модела референтната последователност завършва с Nkx2.2 около 0,94, а вариантът с около 0,008, и връщането на базата възстановява референтната траектория точно. Втори компонент е предсказвател с Huber-линеен модел за ефекта на редакция върху репортерната активност: при дадена последователност или таблица с редакции той връща знакови оценки, подрежда редакциите и избира панел с по една редакция от семейство. Върху 16\,986 вътрешни редакции той намалява средноквадратичната грешка, претеглена по семейства, с 3,02\% спрямо основен модел с 5-мери. Доказателството е предварително: разработката е ползвала тези данни, а върху пет външни изследвания корелацията е 0,049 и грешката е с 9,52\% над тази при предсказване на нулев ефект. Epigenesis Studio е локално настолно приложение на английски и български, което изпълнява симулацията и предсказвателя и изнася удостоверения, които всеки със същата инсталация може да възпроизведе.

\smallskip\noindent\textbf{Ключови думи:} ДНК компилатор, регулаторни варианти, симулация, невронно развитие, възпроизводимост, предсказване на ефекта на варианти.
\end{abstractbox}

\section{Увод}

Изследовател, който променя регулаторен ДНК сайт, иска да знае какво прави промяната с една развиваща се система, по възможност преди да избере лабораторен експеримент. За отговора е нужна верига от стъпки: кои бази се променят, как промяната влияе на свързването, как това влияние влиза в динамичен модел и как моделът реагира във времето. На практика веригата живее в скриптове, които не записват допусканията си. Затова резултатът не може да се проследи до байтовете на източника и друг човек не може да го повтори.

Epigenesis третира ДНК записа като изходен код на програма. Профилът на източника определя граматиката и идентичността на записа. Потребителят задава тълкуване, което казва какво се приема, че означава ДНК. Компилаторът свежда двете до типизиран модул за развитие, независим валидатор повтаря извеждането, а изпълнимите модели четат запечатания модул. Компилацията се приема само ако тази втора реализация я възпроизведе.

Докладът следва реда, в който изследователят среща софтуера. Раздел~\ref{sec:software} описва компилатора, валидатора и изпълнимите експерименти. Раздел~\ref{sec:example} е основният работен процес: промяна на ДНК сайт, симулация на невронна регулаторна мрежа при зададени условия и сравняване на варианта с референтната последователност и с връщането на промяната. Раздел~\ref{sec:predictor} описва предсказвателя на редакции, втори компонент, който оценява и подрежда редакции, с данните му, метода, резултатите, демонстрация в Studio и границите на преноса му. Раздел~\ref{sec:studio} описва Epigenesis Studio, което изпълнява и двата процеса, а раздел~\ref{sec:limits} посочва границите.

Сред разработките, разгледани в раздел~\ref{sec:related}, няма такава, която да компилира съществуваща ДНК при изрично тълкуване и с независимо повторно изпълнение. Това е ограничено сравнение; не сме правили пряк тест срещу други инструменти. Компилаторът не извежда биология от последователността, защото тълкуването е входна величина. Симулациите важат при заявените си допускания и не предсказват изхода при ембриони.

\section{Сходни разработки}\label{sec:related}

\textbf{Компилатори за ДНК.} Няколко системи компилират описание на проект до ДНК или генетични вериги: атрибутни граматики \citep{cai2009}, пространствен език, компилиран до регулаторни мрежи \citep{beal2011}, езиците Eugene и Genotype Specification \citep{bilitchenko2011,wilson2016} и Cello \citep{nielsen2016}. Epigenesis работи в обратната посока: взема съществуваща ДНК и изрично тълкуване и ги свежда до вход за модел.

\textbf{Симулирано развитие.} CX3D израства неврони и мрежи от локални правила \citep{zubler2009}. Управлява го език за инструкции \citep{zubler2011}, приложен и към развитието на кората \citep{zubler2013}. Един от експериментите тук използва CX3D.

\textbf{Възпроизводими модели.} Архивите COMBINE пакетират модела заедно с данните и метаданните му \citep{bergmann2014}. Epigenesis е по-тесен. Той определя какво трябва да записва извеждането от ДНК.

\textbf{Биологични модели в експериментите.} Отчитането на градиента на Sonic Hedgehog в невралната тръба е добре количествено описан модел на решенията за клетъчна съдба \citep{balaskas2012,cohen2014,herreradelgado2020}. Данни за енхансера дават \citet{peterson2012}, а за афинитета на свързване \citet{hallikas2006}.

\textbf{Предсказване на ефекта на варианти.} Раздел~\ref{sec:predictor} и приложение~\ref{app:record} използват данни от репортерни анализи \citep{melnikov2012,tewhey2016,kircher2019} и сравнения с модели от последователност към функция \citep{zhou2015,kelley2016,kelley2018,avsec2021,chen2022}. Предсказвателят съчетава регресия на Хюбер \citep{huber1964} с мотивни оценки от JASPAR \citep{rauluseviciute2024}. И двата метода са утвърдени. Нашият принос е реализираният и оценен работен процес около тях.

\section{Софтуерът}\label{sec:software}

\subsection{Компилатор и договори}

Компилаторът превръща ДНК източник в модул за развитие на четири етапа:
\begin{center}ДНК източник $\rightarrow$ типизиран пакет на източника $\rightarrow$ изрично тълкуване $\rightarrow$ модул за развитие\end{center}
Вградените профили четат суров IUPAC, FASTA, поточен референтен FASTA, GenBank и GFF3. Други граматики влизат през външен интерфейс от версия~2. Той е изпълним файл, закрепен по хеш, със самостоятелен изпълним валидатор. Потребителят задава тълкуването (манифест, заявка, отговор), политиката за свеждане и целевия договор. Компилаторът проверява типовете, координатите, препратките и ограниченията по ресурси.

Всеки артефакт е затворен JSON обект. На всяко ниво компилаторът отхвърля непознати или липсващи членове, дублирани ключове, некрайни числа, цели числа извън ±(2\textsuperscript{53}−1) и грешни типове. Идентичностите са SHA-256 хешове на каноничното кодиране по RFC~8785 \citep{rundgren2020,nist2015}. Идентичностите на последователностите следват схемата refget на GA4GH. \code{SPEC.md} описва договорите. \code{docs/reference/STANDARDS.md} свързва всеки договор със съответната спецификация.

\subsection{Независима проверка}

Валидатор, изолиран от производителя на артефакта, чете наново оригиналния източник, изпълнява външните валидатори върху нови моментни копия и повтаря свеждането и свързването. Той не внася кода на производителя. Проверката при издаване изпълнява 505 модулни теста (на тази машина три се пропускат) и цялостни проверки на минималния пример, универсалния превод, граничните условия и инсталирането от wheel и от архив с изходен код. Минималният пример възпроизвежда четири фиксирани хеша: на описанието на източника, на модула за развитие, на пакета за развитие и на отчета от валидацията.

\subsection{Изпълними експерименти}\label{sec:experiments}

Всеки експеримент посочва входните данни, допусканията, контролите и твърдението, което подкрепя (таблица~\ref{tab:experiments}). Оформянето на невралната тръба е работният пример от раздел~\ref{sec:example}.

\begin{table}[tb]
\small
\caption{Изпълними експерименти.}\label{tab:experiments}
\begin{tabular}{@{}L{25mm}L{56mm}L{67mm}@{}}
\toprule
Експеримент & Какво прави & Резултат и допускане \\
\midrule
Оформяне на невралната тръба & Чете заместване на енхансер на \emph{Nkx2.2} \citep{peterson2012}, измерени съотношения на свързване на GLI \citep{hallikas2006} и модел с четири гена \citep{herreradelgado2020,cohen2014} & Връщането на заместването възстановява изходното състояние. Контролите за история на сигнала, спасяване и фенокопие и числовото уточняване минават. Съотношението на свързване in vitro се пренася към вътреклетъчния афинитет по допускане. \\
Кодиращи последствия & Възстановява промени в \emph{Sox2} от геномен FASTA и ги превежда със стандартния генетичен код на NCBI & Делеция на четири бази дава продукт от 97 аминокиселинни остатъка (69 непроменени) и стоп при кодон 98, което съвпада с мутанта при \citet{gong2020}. Не се твърди нищо за количеството на протеина. \\
Реконструкция на измервания (neural evidence) & Възстановява 109 наблюдения на експресията на \emph{gcy-22} след изчерпване на CHE-1 при \emph{C. elegans} \citep{traets2021}, всяко проследено до клетка от работната книга и до замислена промяна в ДНК & Средните съотношения ауксин/контрола са 8,70\% при пренесено ядро за свързване от 12 бази и 66,33\% при по-голям фрагмент. Ограниченията на извадката и противоречивите етикети в източника остават прикрепени към данните. \\
Регулаторна динамика (neural fate) & Изчислява динамиката на Pax6, Olig2 и Nkx2.2 от изкуствена рецепта с параметри от 512 бази & Възпроизвежда зададено поведение на публикувани модели с обикновени диференциални уравнения. Уравненията задават взаимодействията. Рецептата не ги открива. \\
Изграждане на нервна мрежа & Израства клетки, неврити и контакти от една основаваща клетка и изкуствена рецепта от 480 бази (15 параметъра) върху CX3D; стимулира клетките и прекъсва контакти & Морфологията и връзките възникват по време на изпълнението и не са кодирани. Рецептата е изкуствена. \\
\bottomrule
\end{tabular}
\end{table}

Два по-ранни експеримента завършиха с отрицателен резултат и остават в хранилището (приложение~\ref{app:record}). Експериментите показват, че веригата от ДНК през изрично тълкуване до механистичен модел може да се изпълни и да се провери. Те не показват, че тълкуванията са верни.

\section{Работен пример: промяна на сайт за GLI}\label{sec:example}

\subsection{Въпросът и веригата}

Въпросът е как зададена промяна в ДНК изменя динамиката на модел на невронно развитие. Примерът е експериментът \code{examples/neural-patterning}. Всяка стъпка по-долу е отделна част от данните или кода в хранилището и всяко допускане е заявено.

\begin{enumerate}
\item \textbf{Сайт.} \citet{peterson2012} тестват замяна в сайт за свързване на GLI в енхансера на \emph{Nkx2.2} на $-2$~kb. В ориентацията на веригата от таблицата за свързване деветбазовото ядро се променя от GACCACCCA на GAACACCCA: замяна на C с A в позиция~3. Хранилището съдържа само тези ядра; пълният енхансер не е включен.
\item \textbf{Свързване.} \citet{hallikas2006} измерват свързването на GLI2 и GLI3 с всяка замяна на една база в консенсусния сайт. Съотношението $q$ между мутантния и референтния сайт за тази замяна е 0,0424 за GLI2 (основният избор) и 0,0353 за GLI3 (проверка на чувствителността). Компилаторът превръща избрания сайт в модул за развитие, който съхранява само избраните нуклеотиди (тензорите \code{count} и \code{nucleotides}). Симулацията чете тези бази обратно от валидирания модул и извежда $q$ от таблицата на свързването.
\item \textbf{Пренос.} По заявено допускане $q$ умножава афинитета към GLI на промотора на \emph{Nkx2.2} в модела с четири гена на \citet{herreradelgado2020}, чийто механизъм на конкурентно свързване следва \citet{cohen2014}. Нека $S$ е делът на GLI в активаторна форма, $w$ е базовото тегло, $K$ е коефициентът за GLI, $c$ е коефициентът на активиране и $D$ е факторът на кръстосана репресия. Промоторът е активен с вероятност $P$, тоест неговата заетост:
\[
P=\frac{w(1+K)(1+qcKS)}{w(1+K)(1+qcKS)+(1+qK)D},\qquad \frac{dX}{dt}=2P-2X,
\]
където $X$ е концентрацията на генния продукт в единици на модела. За референтния сайт $q=1$ и факторът $1+K$ се съкращава, което възстановява публикуваните уравнения. При фиксирано ендогенно състояние по-слабият сайт намалява заетостта при сигнали над \mbox{$S=1/c=0{,}1$} и я увеличава под тази стойност.
\item \textbf{Динамика.} Четирите гена Pax6, Olig2, Nkx2.2 и Irx3 се интегрират с адаптивни стъпки по метода на Рунге–Кута от четвърти ред при историята на сигнала и началното състояние, които задава потребителят. Концентрациите и времето са в единици на модела.
\end{enumerate}

\subsection{Референтна последователност, вариант и връщане}

Изпълнихме референтния сайт и варианта при един протокол: дял на GLI активатора 0,2 за 20 единици на модела, като и четирите гена започват от нула. Вариантът се изпълнява по два начина. Като редакция вътре в мрежата (ендогенен режим) променената стойност на афинитета се връща обратно във всичките четири гена. Като пасивен репортер (режим на репортер) промоторът отчита заетост върху непроменената референтна мрежа. Третото изпълнение връща базата и повтаря референтното.

\begin{table}[tb]
\small
\caption{Работен пример при 20 единици на модела: дял на GLI активатора 0,2, профил на GLI2, всички гени започват от нула. Стойностите са концентрации в единици на модела; заетостта е вероятността промоторът на \emph{Nkx2.2} да е активен.}\label{tab:example}
\begin{tabular}{@{}L{54mm}rrrrr@{}}
\toprule
Изпълнение & Pax6 & Olig2 & Nkx2.2 & Irx3 & Заетост \\
\midrule
Референтна последователност & 0,00519 & 0,0208 & 0,940 & 0,000716 & 0,940 \\
Вариант, редакция вътре в мрежата & 0,241 & 0,967 & 0,00846 & 0,00206 & 0,00846 \\
Вариант, пасивен репортер & 0,00519 & 0,0208 & 0,940 & 0,000716 & 0,928 \\
Връщане на базата & \multicolumn{5}{l}{идентично с референтното; най-голяма абсолютна разлика 0} \\
\bottomrule
\end{tabular}
\end{table}

\begin{figure}[tb]
\centering
\patterningplot{време на модела}{референтна последователност}{вариант, редакция вътре в мрежата}{вариант като пасивен репортер}{Заетост на промотора на Nkx2.2}
\caption{Траекториите зад таблица~\ref{tab:example}. В четирите панела за гени пунктираната линия е вариантът като редакция вътре в мрежата. В панела за заетостта точковата линия е вариантът като пасивен репортер, който остава близо до референтната (в края 0,928 срещу 0,940). Режимът на репортер оставя и четирите гена върху референтната траектория.}\label{fig:trajectories}
\end{figure}

С по-слабия сайт Nkx2.2 вътре в мрежата завършва с 0,0085, докато референтната последователност завършва с 0,940. Останалите гени следват. В този модел Nkx2.2 потиска Pax6, Olig2 и Irx3, затова по-малкото Nkx2.2 ги освобождава: Olig2 расте от 0,021 до 0,967, а Pax6 от 0,0052 до 0,241. При отчитане през пасивен репортер същата замяна променя заетостта на промотора само от 0,940 на 0,928 и оставя и четирите гена върху референтната траектория. Репортерен анализ би показал малка промяна при редакция, чийто ефект вътре в мрежата с обратна връзка е голям. При GLI3 вместо GLI2 вариантът завършва с Nkx2.2 0,0081.

Връщането на базата дава траектория, идентична с референтната във всичките 201 момента на изхода и за четирите гена (най-голяма абсолютна разлика 0), а върнатият сайт се компилира до същия модул за развитие като референтния. Програмата на примера добавя дванайсет проверки на механизма и всички минават при демонстрационното изпълнение. Спасяването чрез медиатор запазва ДНК на варианта и налага референтния множител. Фенокопието запазва референтната ДНК и налага множителя на варианта. Контролът за история на сигнала показва, че еднакви крайни сигнали могат да оставят различни състояния на Nkx2.2. Други проверки обхващат обръщането на знака на ефекта от свързването при различни сигнали и сходимостта, когато интервалът на изхода се преполови (разлики под $10^{-4}$).

\subsection{Допускания и обхват}\label{sec:assumptions}

Резултатът зависи от шест заявени допускания, изброени в \code{examples/neural-patterning/data.json} и показвани при всяко изпълнение в Studio.
\begin{enumerate}
\item Публикуваното нормализирано съотношение от профила на свързване се пренася като относителен множител на афинитета. Първоначалният процес на нормализиране не е напълно възстановен.
\item Един и същ GLI профил мащабира афинитета на активатора и на репресора. GLI2 е основният избор, а GLI3 е анализ на чувствителността.
\item Деветбазовият сайт е отнесен към публикуван регулаторен контекст на \emph{Nkx2.2}. Останалата последователност на енхансера, хроматинът и кофакторите са фиксирани и не се моделират.
\item Режимът на репортер отчита заетостта на промотора върху непроменена ендогенна траектория. Ендогенният режим е хипотетична редакция на генома с обратна връзка.
\item Сигналът е делът на GLI активатора в $[0,1]$. Той не е калибрирана концентрация на Shh.
\item Времето и концентрацията са в единици на модела. Не се извежда флуоресцентна скала, праг на откриване, биологичен часовник или вероятност за ембрион.
\end{enumerate}
При \citet{peterson2012} репортерът на \emph{Nkx2.2} е експресиран при 4 от 6 ембриона с референтния сайт и при 0 от 13 с варианта (E10.5). Моделът няма праг на откриване и модел на шума на ниво ембрион, затова тези броеве се показват до симулацията и не се оценяват като възпроизведено предсказание. Моделът с четири гена не включва Foxa2.

Поддържаният вход е една замяна с измерено положително тегло на свързване в този сайт, протокол на сигнала до 10 фази и 100 единици на модела и по избор начално състояние. Друг геном или сайт е извън поддържания път. Резултатът е условна механистична симулация. Тя помага да се сравняват хипотези за промяна на свързването, преди да се избере експеримент, и не доказва развитие от естествен геном до мозък, свързаност или учене.

\section{Предсказвателят на редакции}\label{sec:predictor}

Симулацията от раздел~\ref{sec:example} следва една избрана промяна през механистичен модел. Предсказвателят разглежда много кандидат-редакции наведнъж. За всяка той оценява само по последователността с колко редакцията променя активността на репортера и в каква посока. Изследователят може да ползва оценките, за да подреди редакциите и да избере панел за измерване. Studio изпълнява предсказвателя в разделите Сканиране и Редакции (раздел~\ref{sec:demo}).

\subsection{Задача и данни}\label{sec:data}

\textbf{Задача.} При дадени референтна и редактирана последователност да се оцени ефектът върху активността на репортера, алтернативен минус референтен ($\log_2$ на отношението), в 24-часовия анализ в SK-N-SH на \citet{siraj2026}. Допускат се замени, вмъквания и изтривания. Оценката подрежда редакциите в рамките на този анализ. Тя не е вероятност и не симулира развитие.

\textbf{Скрининг и дялове.} Скринингът на Siraj има 365\,762 двойки конструкти, 336\,475 от които с два алела по 200 бази, измерени с поне пет повторни трансфекции на библиотека. Редакциите са групирани в семейства, свързани компоненти от припокриващи се или идентични конструкти (включително обратно-комплементарни), така че сродни редакции никога не попадат от двете страни на разделянето. Крайният модел е напасван върху дяла за напасване (FIT), 66\,604 редакции в 6794 семейства, и е сравняван върху вътрешния дял за настройка (TUNING), 16\,986 редакции в 1759 семейства. Още един дял, RESERVE, не е оценяван и остава запечатан. Разработката е ползвала TUNING многократно, затова резултатите върху него са резултати от настройване (приложение~\ref{app:protocol} изброява четирите класа данни).

\textbf{Външни изследвания.} Пет невронни репортерни изследвания проверяват преноса (таблица~\ref{tab:transfer}). Отделеният набор на всяко от тях съдържа само варианти, при които няма вариант от друго изследване и няма конструкт на Siraj в рамките на 2000 бази, и е един и същ за всички модели. По-ранни модели на проекта вече бяха виждали тези изследвания, затова данните са разкрити. Шесто изследване, Rummel и сътр., е анализ на чувствителността, защото при 83\% от вариантите единият алел има по-малко от десет RNA прочита (приложение~\ref{app:record}).

\subsection{Метод}\label{sec:method}

\textbf{Признаци.} За една редакция моделът образува разликата между редактираната и референтната последователност в два блока признаци (фигура~\ref{fig:pipeline}). Първият брои 512-те канонични 5-мера, като 5-мер и обратният му комплемент се броят заедно. Вторият има 532 мотивни признака. За всеки мотив на гръбначни от JASPAR 2024 \citep{rauluseviciute2024} потенциалът на една последователност е сумата по прозорците ѝ от $\log(1+s\,r)$, където $r$ е шансът на мотива в прозореца, усреднен по двете вериги, а $s=0{,}003$ компресира силните сайтове. Потенциалът се мери спрямо прозорец с $r=1$. Признакът е потенциалът на редактираната последователност минус този на референтната. Различават се само прозорците, които припокриват редакцията, затова оценката зависи от 18-те бази от двете страни на промяната и от нищо по-далечно.

\begin{figure}[htb]
\centering
\pipelinefig{референтна и редактирана последователност}{18 бази от двете страни на промяната}{512 разлики в броя на 5-мерите}{532 разлики в потенциалите на мотивите}{знаков ефект от линеен модел, напасван с загуба на Хюбер}{подредба по абсолютен ефект, една редакция на семейство: панелът}
\caption{Предсказвателят. Редакцията променя два блока признаци. Линеен модел ги превръща в знаков ефект, а панелът взема най-силните редакции, по една от семейство. Studio изпълнява тази контролна точка.}\label{fig:pipeline}
\end{figure}

\textbf{Напасване.} Двата блока се напасват заедно, без свободен член, като се минимизира загуба на Хюбер \citep{huber1964} с наказание 0,01 за блока с 5-мери и 3 за блока с мотиви и с равно тегло на всяко семейство. Прагът на Хюбер се задава веднъж, от медианното абсолютно отклонение на остатъците от напасването с квадратична загуба, и не се променя по време на оптимизацията. Всяко напасване се е сходило за 10 до 13 актуализации. Контролната точка е \code{epigenesis-siraj-huber-linear-v1}.

\textbf{Базови модели.} \emph{k5} е гребенова регресия само върху блока с 5-мери. \emph{Joint-linear} напасва двата блока със същите наказания и квадратична загуба, така че сравнението му с Huber-линейния модел изолира загубата. \emph{Без ефект} предсказва, че никоя редакция не променя нищо. \emph{S2}, който се ползва в раздел~\ref{sec:transfer}, е по-ранният модел, обучен върху Siraj, който добавя потенциали на мотиви към k5 (приложение~\ref{app:record}).

\textbf{Панели и полезност.} Панел от $n$ редакции се строи, като редакциите се подредят по абсолютен предсказан ефект (равенствата по идентификатор) и се взема първата редакция от всяко още непредставено семейство, докато се съберат $n$. Знаковата полезност на панела е средният наблюдаван ефект на редакциите му, всеки умножен по знака на предсказания си ефект. При случаен подбор стойността е около нула. Интервалите идват от 2000 преизвадки на цели семейства, а долната граница е 5-ият персентил на печалбата спрямо сравнявания модел.

\subsection{Резултати в рамките на анализа}\label{sec:internal}

Таблица~\ref{tab:internal} сравнява модела с базовите му модели във вътрешния дял за настройка. Huber-линейният модел има най-ниската грешка: с 3,02\% под k5 и с 5,40\% под предсказването на нулев ефект. Joint-linear се различава от него само по загубата, а Huber-линейният го надминава с 0,0030 при положителна долна граница (+0,0002). При панели от 50 и 200 редакции Huber-линейният модел има най-висока полезност (фигура~\ref{fig:utility}). При десет редакции joint-linear е по-горе (1,601 срещу 1,483), защото двата панела се различават в две редакции, и данните не подреждат моделите там.

\begin{table}[htb]
\small
\caption{Вътрешен дял за настройка, 16\,986 редакции в 1759 семейства. Грешката е средноквадратичната грешка, претеглена по семейства (по-ниска е по-добре). Полезността е знаковата полезност на панел от 10, 50 или 200 редакции (по-висока е по-добре). Последният ред е долният край на интервала от 5-ия до 95-ия персентил на печалбата на Huber-linear спрямо k5 при 2000 преизвадки на цели семейства.}\label{tab:internal}
\begin{tabular}{@{}L{50mm}rrrr@{}}
\toprule
Модел & Грешка & Полезност, 10 & Полезност, 50 & Полезност, 200 \\
\midrule
Без ефект & 0,304075 & 0 & 0 & 0 \\
k5 & 0,296611 & 0,857 & 0,607 & 0,509 \\
Joint-linear (квадратична загуба) & 0,290683 & 1,601 & 0,954 & 0,626 \\
Huber-linear (краен) & 0,287667 & 1,483 & 0,986 & 0,634 \\
\midrule
Печалба спрямо k5, долна граница & +0,0051 & −0,150 & +0,140 & +0,063 \\
\bottomrule
\end{tabular}
\end{table}

\begin{figure}[htb]
\centering
\utilityplot{k5}{Joint-linear}{Huber-linear}{Размер на панела (редакции)}{Знакова полезност}
\caption{Знакова полезност на панелите по размер във вътрешния дял за настройка (стойностите са в таблица~\ref{tab:internal}).}\label{fig:utility}
\end{figure}

Шест от седемте първоначални изисквания са изпълнени. Отвореното е печалбата при десет редакции спрямо k5. Точковата ѝ оценка е положителна (+0,626), а долната ѝ граница е отрицателна (−0,150), така че данните не го разрешават. Разработката е ползвала този дял през цялото време, затова наричаме подобрението предварително.

\subsection{Демонстрация в Studio}\label{sec:demo}

Studio изпълнява крайната контролна точка. Демонстрацията ползва два входа, които се доставят с приложението: последователност от 200 бази и таблица от 64 редакции в 11 семейства без измерени стойности, така че всяко число по-долу е предсказание. Скриптът \code{docs/paper/figure-data.py} възпроизвежда всяка стойност.

\begin{figure}[htb]
\centering
\scanplot{Позиция в последователността от 200 бази}{Предсказан ефект}{G към A в 100: $-0{,}214$}
\caption{Сканиране в Studio на примерната последователност. Всяка колона е позиция, а всеки ред е базата, в която се променя; празната клетка в колона е референтната база. Оранжевите клетки са предсказани увеличения, а сините са намаления на репортерната активност, наситени при $\pm0{,}6$. Стрелката сочи първата редакция на примерната таблица.}\label{fig:scan}
\end{figure}

\textbf{Сканиране.} Поставянето на последователността от 200 бази в раздела Сканиране оценява всичките 600 замени на една база (фигура~\ref{fig:scan}). От 600-те редакции 343 се предсказват да повишават репортерната активност, а 257 да я понижават, и медианният абсолютен ефект е 0,049. Най-голямото повишение е A към G в позиция 52 (+0,581), а най-голямото понижение е G към T в позиция 165 (−0,313). Първата редакция на примерната таблица, G към A в позиция 100, получава −0,214 при Сканиране. Разделът Редакции връща същата стойност за същата редакция като двойка, защото оценката зависи само от базите около промяната.

\textbf{Редакции.} Разделът Редакции приема двойки референтна и редактирана последователност, всяка с по избор семейство, и размер на панела. За панел от пет върху таблицата с 64 редакции петте най-големи абсолютни оценки включват две редакции от едно семейство (локусът chr10:102622941, +0,169 и +0,167). Правилото за подбор запазва първата и преминава към следващото непредставено семейство, така че панелът има пет редакции от пет семейства (таблица~\ref{tab:panel}). Всяка оценка е сума от част за 5-мерите и част за мотивите, които удостоверението записва поотделно.

\begin{table}[htb]
\small
\caption{Панел от пет от примерната таблица с 64 редакции. Семейство е етикетът, който таблицата дава на всяка редакция. Ефектът е предсказаният ефект, алтернативен минус референтен, сумата от частта за 5-мерите и частта за мотивите.}\label{tab:panel}
\begin{tabular}{@{}r L{36mm} L{22mm} r r r L{16mm}@{}}
\toprule
№ & Редакция & Семейство & Ефект & Част 5-мери & Част мотиви & Посока \\
\midrule
1 & chr10:101119242 G>A & 101119242 & −0,2140 & −0,1261 & −0,0878 & понижава \\
2 & chr10:119062236 C>T & 118615567 & −0,2033 & −0,1434 & −0,0599 & понижава \\
3 & chr10:102798346 T>C & 102622941 & +0,1687 & +0,1407 & +0,0280 & повишава \\
4 & chr10:105307889 T>C & 105307889 & −0,1572 & −0,1439 & −0,0133 & понижава \\
5 & chr10:119050998 C>T & 119021407 & +0,1326 & +0,0162 & +0,1164 & повишава \\
\bottomrule
\end{tabular}
\end{table}

\textbf{Удостоверение.} Изнасянето на резултата от Редакции записва удостоверение. Зареждането му в раздела Проверка или командата \code{epigenesis-studio verify} повтаря заявката и преминава и трите проверки, с максимална разлика 0 между записаните и повторно изчислените числа. Таблица~\ref{tab:verify} показва променени копия: променен резултат пропада на повторното изчисляване дори при преизчислена контролна сума, а преименуван модел пропада на проверката за идентичност. Удостоверението казва на читателя, че резултатът идва от тази контролна точка и тази заявка. То не казва нищо дали предсказанието е вярно.

\begin{table}[htb]
\small
\caption{Проверка на удостоверението на панела по-горе и на три променени копия. Последната колона дава най-голямата разлика между записаните и повторно изчислените числа.}\label{tab:verify}
\begin{tabular}{@{}L{62mm}L{20mm}L{26mm}L{26mm}@{}}
\toprule
Удостоверение & Контролна сума & Идентичност на модела & Повторно изчисляване \\
\midrule
Изнесено & да & да & да (0) \\
Първият записан ефект повишен с 0,001 & не & да & не (0,001) \\
Същото, контролната сума преизчислена & да & да & не (0,001) \\
Модел преименуван, контролната сума преизчислена & да & не & да (0) \\
\bottomrule
\end{tabular}
\end{table}

\subsection{Пренос към други изследвания}\label{sec:transfer}

Същата контролна точка беше оценена веднъж, без ново напасване, върху петте невронни репортерни изследвания (таблица~\ref{tab:transfer}). Корелациите по изследвания са малки и с различен знак. Средната е 0,0490 срещу 0,0465 за S2 и 0,0460 за k5. Сдвоената печалба спрямо S2 е +0,0024 с интервал от −0,0060 до +0,0109, така че корелациите не разделят моделите.

Грешката е по-строгият тест. Huber-линейният модел има по-ниска средноквадратична грешка от всеки напасван сравнителен модел. Средната му стойност 0,1044 обаче е с 9,52\% над 0,0953 на предсказвател, който казва, че никоя редакция няма ефект, и всяка от петте точки по изследвания и техните интервали е в полза на липсата на ефект (фигура~\ref{fig:transfer}). В рамките на анализа същият модел е с 5,40\% под този предсказвател. Резултатите в изследванията се различават по определение и мащаб, а сравнението с липса на ефект беше обявено след оценяването, така че това е диагностика, а не повторение на експеримента.

Оценките подреждат редакциите малко по-добре от случайността в тези изследвания, а мащабът им не се пренася от един анализ към друг. Не твърдим пренос към други анализи, типове клетки или ендогенна експресия. Тринайсет фиксирани теста (приложение~\ref{app:tests}) не намериха модел, който надеждно да надмине простите базови модели в различни контексти.

\begin{table}[htb]
\footnotesize
\caption{Huber-linear върху отделените редакции на петте невронни репортерни изследвания, оценен веднъж без ново напасване. $r$ е корелацията на Пиърсън, претеглена по семейства. Грешката е средноквадратичната грешка на Huber-linear и на предсказването на нулев ефект, а Излишъкът е първата като процент над втората. Редовете със средни стойности дават на изследванията равно тегло; излишъкът в реда за средното сравнява двете средни грешки. Matoba има 97 семейства. $^\dagger$Изследване за чувствителност, извън средното.}\label{tab:transfer}
\begin{tabular}{@{}L{25mm}rrrrrrrrr@{}}
\toprule
 & & & \multicolumn{3}{c}{Корелация $r$} & \multicolumn{3}{c}{Грешка} & \\
\cmidrule(lr){4-6}\cmidrule(lr){7-9}
Изследване & Редакции & Семейства & Huber & S2 & k5 & Без ефект & Huber & Излишък \\
\midrule
Lee & 5916 & 2340 & 0,0152 & 0,0185 & 0,0254 & 0,0626 & 0,0739 & +18,10\% \\
McAfee & 1682 & 676 & 0,0935 & 0,0910 & 0,0990 & 0,0277 & 0,0333 & +20,38\% \\
Matoba & 4013 & 97 & 0,0286 & 0,0225 & 0,0384 & 0,0164 & 0,0256 & +56,53\% \\
Guo & 748 & 571 & −0,0157 & −0,0143 & −0,0241 & 0,3479 & 0,3612 & +3,83\% \\
Weiss & 9247 & 6930 & 0,1233 & 0,1149 & 0,0912 & 0,0220 & 0,0278 & +26,46\% \\
\midrule
Средно от петте & & & 0,0490 & 0,0465 & 0,0460 & 0,0953 & 0,1044 & +9,52\% \\
Rummel$^\dagger$ & 611 & 349 & 0,0375 & 0,0488 & 0,0505 & 0,5420 & 0,5503 & +1,53\% \\
\bottomrule
\end{tabular}
\end{table}

\begin{figure}[htb]
\centering
\correlationplot{Корелация $r$}{Huber-linear}{S2}{k5}\hfill\excessplot{Грешка над липсата на ефект (\%)}{Вътрешен}{Средно}
\caption{Вляво: корелация по изследвания (таблица~\ref{tab:transfer}). Вдясно: грешка на Huber-linear спрямо предсказването на нулев ефект. Първата лента е вътрешният дял за настройка (под липсата на ефект), следващите пет са изследванията, а последната сравнява двете средни грешки.}\label{fig:transfer}
\end{figure}

\subsection{Предназначение}

Предсказвателят е предназначен за подреждане на редакции в рамките на изследвания анализ, например при избор коя от много кандидат-редакции да се измери първа. Дали това спестява лабораторна работа, не е измерено. Извън анализа оценката е хипотеза за проверка.

\section{Epigenesis Studio}\label{sec:studio}

Studio (\code{studio/}, команда \code{epigenesis-studio}) е локално настолно приложение. Сървърът му слуша само на 127.0.0.1 и отхвърля заявки, чийто Host заглавен ред не е локален. Интерфейсът е на английски и български с превключвател в страничната лента. Приложението има четири раздела (таблица~\ref{tab:studio}).

\begin{table}[tb]
\small
\caption{Разделите на Studio.}\label{tab:studio}
\begin{tabular}{@{}L{22mm}L{60mm}L{66mm}@{}}
\toprule
Раздел & Вход & Изход \\
\midrule
Симулация & Една замяна в деветбазовия сайт, GLI2 или GLI3, протокол на сигнала до 10 фази и 100 единици на модела и по избор начално състояние & Траектории за референтната последователност, варианта и връщането; сравнение между репортер и мрежа; контроли; шестте допускания; CSV; удостоверение \\
Сканиране & Един FASTA запис или обикновена последователност от 5 до 1000 бази (само A, C, G, T; няколко записа се отказват) & Топлинна карта на трите замени във всяка позиция; подредена таблица; CSV; удостоверение \\
Редакции & До 5000 двойки референтна и редактирана последователност (CSV или JSON) с по избор семейство за всяка & Оценки; панел с поискания размер с по една редакция от семейство; CSV; удостоверение \\
Проверка & Файл с удостоверение & Три проверки: контролна сума, идентичност на модела и крайната точка, повторно изчисляване на всеки записан резултат \\
\bottomrule
\end{tabular}
\end{table}

Симулация изпълнява програмата на примера от раздел~\ref{sec:example} на живо; втори симулатор няма. Тя изисква копие на хранилището; без него разделът е изключен, а останалите работят.

Сканиране и Редакции изпълняват предсказвателя от раздел~\ref{sec:predictor}, който раздел~\ref{sec:demo} демонстрира. Сканиране оценява всяка промяна на една база в прозорец около позицията, което стига, защото оценката зависи само от 18-те бази от двете страни. Резултатът е равен на оценяването на цели последователности (проверено до $10^{-12}$). Редакции оценява таблица от двойки и избира панела по правилото от раздел~\ref{sec:method}.

Удостоверението (\code{epigenesis-studio-receipt-v1}) съдържа името на модела, хеша на идентичността му, крайната точка, нормализираната заявка, резултата и SHA-256 контролна сума върху каноничния JSON по RFC~8785 на всичко това. Проверка сверява контролната сума, че името, идентичността и крайната точка на модела съвпадат с инсталацията и че повторното изпълнение на записаната заявка възпроизвежда всеки записан резултат с точност $10^{-9}$ (таблица~\ref{tab:verify}).

Само Симулация компилира входа си. Сканиране и Редакции ползват каноничния JSON на компилатора за удостоверенията и не минават през веригата на модула за развитие. Тестовете на Studio (21, успешни) обхващат сканирането спрямо оценяване на цели последователности, панела спрямо замразения модел, всяка проверка на удостоверението, ограниченията и отказите при входа, HTTP защитите и симулацията спрямо командния ред на програмата на примера и публикуваното ѝ съотношение на свързване.

\section{Граници и отворени въпроси}\label{sec:limits}

\textbf{Симулация.} Резултатът зависи от шестте допускания от раздел~\ref{sec:assumptions}. Съотношението на свързване идва от измерване in vitro, а използването му като вътреклетъчен множител на афинитета не е проверено. Броевете ембриони на \citet{peterson2012} се показват и не се оценяват. Поддържат се един сайт, един контекст на енхансера и четири гена. Процесът не моделира многоклетъчна механика, свързаност или учене.

\textbf{Предсказвател.} Вътрешните печалби са предварителни, защото разработката е използвала данните, а изискването за десет редакции остава отворено. Раздел~\ref{sec:transfer} дава външната граница. При тринайсет фиксирани теста за пренос (приложение~\ref{app:record}) нито един кандидат не получи етикета си. Точността на посоката за eQTL с фино картиране (от 0,52 до 0,65) съвпадна с тази на базов модел с k-мери, а правилото за най-близкия ген (0,460) избра причинния вариант на достоверния набор по-добре от всеки модел. Затова предсказването между биологични контексти остава отворен проблем. Възможни причини, всички непроверени: репортерните ефекти може да зависят от вектора, промотора и състоянието на клетката; мащабите на ефекта се различават между изследванията; частта от ефектите на единични бази, която се предсказва от последователността, може да е малка; а късите признаци може да улавят главно състав, например съдържание на GC. Предсказване между анализи със същия вектор и тип клетки, на общ мащаб, само с добре измерени варианти и с данни от ендогенни локуси би проверило това.

\textbf{Оценка.} Хешовете на фиксираната процедура са локални. Те пазят от грешки и незаписани промени. Не пазят от автор, който пренаписва плана и хеша му заедно. Тринайсетте теста не са коригирани за множествено тестване. Повторните изчисления са правени от същия екип. Няколко набора за оценка са били разкрити на по-ранни модели. Някои изследвания споделят лабораторна линия или източник на клетки. Данните на Salomon за неврони са изключени от тестовете между изследвания, защото таблиците с варианти са ограничени в портала на IGVF. Крайната точка на предсказвателя е репортерна активност или експресия; развитието не е измерено.

\textbf{Следващи стъпки.} (а) Сравнение на симулираните последствия от примера, например прехода от Nkx2.2 към Olig2, с измерени ендогенни промени и замяна на допуснатите параметри на модела с измерени. (б) Еднократно използване на запечатания дял RESERVE на Siraj за един предварително избран модел. (в) Получаване на ограничените таблици с варианти на IGVF и тест върху тях по нов план. (г) Събиране на набори от данни за ендогенни промени с измерена експресия. Заключението за преноса би се променило, ако модел надмине правилото за най-близкия ген или базов модел с k-мери при запечатан критерий и в два независими набора от данни, които не е виждал.

\section{Заключение}

Epigenesis позволява на изследователя да промени поддържан ДНК сайт, да изпълни модел на развитие при зададени условия и да разгледа последствията, като всяка стъпка е записана и друга инсталация може да я повтори. В работния пример промяна на една база в сайт за GLI премества Nkx2.2 в модела от 0,94 на 0,008, а връщането възстановява референтната траектория точно. Предсказвателят и Studio добавят начин да се разглеждат и подреждат редакции в рамките на един анализ. Доказателството им е предварително, а предсказването между биологични контексти остава отворено.

\section{Наличност и декларации}

\textbf{Код и данни.} \url{https://github.com/dttdrv/epigenesis}. Кодът е под Apache-2.0. Включените данни и компоненти на трети страни запазват лицензите и авторските си бележки. Например JASPAR 2024 е под CC BY 4.0, бележката за 80K-Analysis е в \code{examples/neural-predictor/LICENSE.Salomon}, а включеният предсказвател с механизъм на внимание използва Interlace (вж. \code{NOTICE.md} в същата папка).

\textbf{Изпълнение.} README съдържа командите. Python 3.11 или по-нов и стандартната библиотека изпълняват компилатора и демонстрациите за реконструкция на измервания и за оформяне. Studio изисква NumPy и компилатора: \code{uv tool install --editable studio --with-editable .} поставя \code{epigenesis-studio} в пътя за изпълнение. Предварително обучените извличащи признаци и експериментът за изграждане на нервна мрежа изискват допълнителни инструменти, които документацията им назовава.

\textbf{Данни за фигурата.} \code{docs/paper/figure-patterning.csv} е записан от симулацията в Studio за GAACACCCA с GLI2 (една фаза от 20 единици на модела при сигнал 0,2, всички гени от нула); \code{patterning.py run} дава същите траектории. \code{docs/paper/figure-data.py} възстановява \code{figure-scan.csv} и \code{figure-studies.csv} и отпечатва числата от демонстрацията в раздел~\ref{sec:predictor}. Полезностите в таблица~\ref{tab:internal} идват от \code{RESULTS.md} на фиксираното съвместно сравнение на Huber в дървото за изграждане без контрол на версиите (\code{build/neural-predictor-reviewed-20261005/siraj-huber-joint-20261019/}).

\textbf{Фиксирани тестове.} Плановете, хешовете, скриптовете и резултатите са в \code{research/frozen-tests/} (приложение~\ref{app:tests}). \code{python3 research/check\_freeze.py} преизчислява хешовете. Дялът RESERVE на Siraj остава запечатан. Потвържденията на Salomon и на Kosicki са изразходвани.

\textbf{Данни без журнална статия.} Серии на Gene Expression Omnibus GSE244011, GSE211045, GSE225815, GSE221384, GSE152404 и GSE188264 (\citealt{kreimer2022}); хранилището 80K-Analysis (\url{https://github.com/kircherlab/80K-Analysis}); Interlace (\url{https://github.com/MisulOrg/Interlace}); CX3D 0.03 (\url{https://tilde.ini.uzh.ch/~amw/seco/cx3d/cx3d-0.03.zip}); GENCODE v39 \citep{frankish2021}; пътеките UCSC hg38 и phyloP100way.

\textbf{Отговорност.} Авторът отговаря за съдържанието. Всички числа идват от файловете с резултати, изброени в приложение~\ref{app:tests}, от папките на примерите и от README на проекта. Конфликт на интереси и финансиране: да се допълнят от автора преди изпращане за публикуване.

\begin{thebibliography}{99}
\setlength{\itemsep}{2pt}

\bibitem[Agarwal et~al.(2025)]{agarwal2025} Agarwal V, Inoue F, Schubach M, et al. (2025). Massively parallel characterization of transcriptional regulatory elements. \textit{Nature} 639:411–420. \href{https://doi.org/10.1038/s41586-024-08430-9}{doi:10.1038/s41586-024-08430-9}

\bibitem[Ashuach et~al.(2019)]{ashuach2019} Ashuach T, Fischer DS, Kreimer A, et al. (2019). MPRAnalyze: statistical framework for massively parallel reporter assays. \textit{Genome Biology} 20:183. \href{https://doi.org/10.1186/s13059-019-1787-z}{doi:10.1186/s13059-019-1787-z}

\bibitem[Avsec et~al.(2021)]{avsec2021} Avsec Ž, Agarwal V, Visentin D, et al. (2021). Effective gene expression prediction from sequence by integrating long-range interactions. \textit{Nature Methods} 18:1196–1203. \href{https://doi.org/10.1038/s41592-021-01252-x}{doi:10.1038/s41592-021-01252-x}

\bibitem[Aygün et~al.(2021)]{aygun2021} Aygün N, Elwell AL, Liang D, et al. (2021). Brain-trait-associated variants impact cell-type-specific gene regulation during neurogenesis. \textit{American Journal of Human Genetics} 108:1647–1668. \href{https://doi.org/10.1016/j.ajhg.2021.07.011}{doi:10.1016/j.ajhg.2021.07.011}

\bibitem[Balaskas et~al.(2012)]{balaskas2012} Balaskas N, Ribeiro A, Panovska J, et al. (2012). Gene regulatory logic for reading the Sonic Hedgehog signaling gradient in the vertebrate neural tube. \textit{Cell} 148:273–284. \href{https://doi.org/10.1016/j.cell.2011.10.047}{doi:10.1016/j.cell.2011.10.047}

\bibitem[Beal et~al.(2011)]{beal2011} Beal J, Lu T, Weiss R. (2011). Automatic compilation from high-level biologically-oriented programming language to genetic regulatory networks. \textit{PLoS ONE} 6:e22490. \href{https://doi.org/10.1371/journal.pone.0022490}{doi:10.1371/journal.pone.0022490}

\bibitem[Bergmann et~al.(2014)]{bergmann2014} Bergmann FT, Adams R, Moodie S, et al. (2014). COMBINE archive and OMEX format: one file to share all information to reproduce a modeling project. \textit{BMC Bioinformatics} 15:369. \href{https://doi.org/10.1186/s12859-014-0369-z}{doi:10.1186/s12859-014-0369-z}

\bibitem[Bilitchenko et~al.(2011)]{bilitchenko2011} Bilitchenko L, Liu A, Cheung S, et al. (2011). Eugene – a domain specific language for specifying and constraining synthetic biological parts, devices, and systems. \textit{PLoS ONE} 6:e18882. \href{https://doi.org/10.1371/journal.pone.0018882}{doi:10.1371/journal.pone.0018882}

\bibitem[Cai et~al.(2009)]{cai2009} Cai Y, Lux MW, Adam L, Peccoud J. (2009). Modeling structure-function relationships in synthetic DNA sequences using attribute grammars. \textit{PLoS Computational Biology} 5:e1000529. \href{https://doi.org/10.1371/journal.pcbi.1000529}{doi:10.1371/journal.pcbi.1000529}

\bibitem[Chen et~al.(2022)]{chen2022} Chen KM, Wong AK, Troyanskaya OG, Zhou J. (2022). A sequence-based global map of regulatory activity for deciphering human genetics. \textit{Nature Genetics} 54:940–949. \href{https://doi.org/10.1038/s41588-022-01102-2}{doi:10.1038/s41588-022-01102-2}

\bibitem[Cohen et~al.(2014)]{cohen2014} Cohen M, Page KM, Perez-Carrasco R, Barnes CP, Briscoe J. (2014). A theoretical framework for the regulation of Shh morphogen-controlled gene expression. \textit{Development} 141:3868–3878. \href{https://doi.org/10.1242/dev.112573}{doi:10.1242/dev.112573}

\bibitem[Deng et~al.(2024)]{deng2024} Deng C, Whalen S, Steyert M, et al. (2024). Massively parallel characterization of regulatory elements in the developing human cortex. \textit{Science} 384:eadh0559. \href{https://doi.org/10.1126/science.adh0559}{doi:10.1126/science.adh0559}

\bibitem[Frankish et~al.(2021)]{frankish2021} Frankish A, et al. (2021). GENCODE 2021. \textit{Nucleic Acids Research} 49:D916–D923. \href{https://doi.org/10.1093/nar/gkaa1087}{doi:10.1093/nar/gkaa1087}

\bibitem[Fromer et~al.(2016)]{fromer2016} Fromer M, Roussos P, Sieberts SK, et al. (2016). Gene expression elucidates functional impact of polygenic risk for schizophrenia. \textit{Nature Neuroscience} 19:1442–1453. \href{https://doi.org/10.1038/nn.4399}{doi:10.1038/nn.4399}

\bibitem[Gneiting and Raftery(2007)]{gneiting2007} Gneiting T, Raftery AE. (2007). Strictly proper scoring rules, prediction, and estimation. \textit{Journal of the American Statistical Association} 102:359–378. \href{https://doi.org/10.1198/016214506000001437}{doi:10.1198/016214506000001437}

\bibitem[Gong et~al.(2020)]{gong2020} Gong et al. (2020). The requirement of Sox2 for the spinal cord motor neuron development of zebrafish. \textit{Frontiers in Molecular Neuroscience} 13:34. \href{https://doi.org/10.3389/fnmol.2020.00034}{doi:10.3389/fnmol.2020.00034}

\bibitem[GTEx Consortium(2020)]{gtex2020} GTEx Consortium (2020). The GTEx Consortium atlas of genetic regulatory effects across human tissues. \textit{Science} 369:1318–1330. \href{https://doi.org/10.1126/science.aaz1776}{doi:10.1126/science.aaz1776}

\bibitem[Guo et~al.(2023)]{guo2023} Guo et al. (2023). Integrative analyses highlight functional regulatory variants associated with neuropsychiatric diseases. \textit{Nature Genetics} 55:1876–1891. \href{https://doi.org/10.1038/s41588-023-01533-5}{doi:10.1038/s41588-023-01533-5}

\bibitem[Hallikas et~al.(2006)]{hallikas2006} Hallikas O, Palin K, Sinjushina N, et al. (2006). Genome-wide prediction of mammalian enhancers based on analysis of transcription-factor binding affinity. \textit{Cell} 124:47–59. \href{https://doi.org/10.1016/j.cell.2005.10.042}{doi:10.1016/j.cell.2005.10.042}

\bibitem[Herrera-Delgado et~al.(2020)]{herreradelgado2020} Herrera-Delgado E, Briscoe J, Sollich P. (2020). Tractable nonlinear memory functions as a tool to capture and explain dynamical behaviors. \textit{Physical Review Research} 2:043069. \href{https://doi.org/10.1103/PhysRevResearch.2.043069}{doi:10.1103/PhysRevResearch.2.043069}

\bibitem[Huber(1964)]{huber1964} Huber PJ. (1964). Robust estimation of a location parameter. \textit{Annals of Mathematical Statistics} 35:73–101. \href{https://doi.org/10.1214/aoms/1177703732}{doi:10.1214/aoms/1177703732}

\bibitem[Jaffe et~al.(2018)]{jaffe2018} Jaffe AE, Straub RE, Shin JH, et al. (2018). Developmental and genetic regulation of the human cortex transcriptome illuminate schizophrenia pathogenesis. \textit{Nature Neuroscience} 21:1117–1125. \href{https://doi.org/10.1038/s41593-018-0197-y}{doi:10.1038/s41593-018-0197-y}

\bibitem[Jerber et~al.(2021)]{jerber2021} Jerber J, Seaton DD, Cuomo ASE, et al. (2021). Population-scale single-cell RNA-seq profiling across dopaminergic neuron differentiation. \textit{Nature Genetics} 53:304–312. \href{https://doi.org/10.1038/s41588-021-00801-6}{doi:10.1038/s41588-021-00801-6}

\bibitem[Kang et~al.(2011)]{kang2011} Kang HJ, Kawasawa YI, Cheng F, et al. (2011). Spatio-temporal transcriptome of the human brain. \textit{Nature} 478:483–489. \href{https://doi.org/10.1038/nature10523}{doi:10.1038/nature10523}

\bibitem[Kelley et~al.(2016)]{kelley2016} Kelley DR, Snoek J, Rinn JL. (2016). Basset: learning the regulatory code of the accessible genome with deep convolutional neural networks. \textit{Genome Research} 26:990–999. \href{https://doi.org/10.1101/gr.200535.115}{doi:10.1101/gr.200535.115}

\bibitem[Kelley et~al.(2018)]{kelley2018} Kelley DR, Reshef YA, Bileschi M, et al. (2018). Sequential regulatory activity prediction across chromosomes with convolutional neural networks. \textit{Genome Research} 28:739–750. \href{https://doi.org/10.1101/gr.227819.117}{doi:10.1101/gr.227819.117}

\bibitem[Kerimov et~al.(2021)]{kerimov2021} Kerimov N, Hayhurst JD, Peikova K, et al. (2021). A compendium of uniformly processed human gene expression and splicing quantitative trait loci. \textit{Nature Genetics} 53:1290–1299. \href{https://doi.org/10.1038/s41588-021-00924-w}{doi:10.1038/s41588-021-00924-w}

\bibitem[Kircher et~al.(2019)]{kircher2019} Kircher M, Xiong C, Martin B, et al. (2019). Saturation mutagenesis of twenty disease-associated regulatory elements at single base-pair resolution. \textit{Nature Communications} 10:3583. \href{https://doi.org/10.1038/s41467-019-11526-w}{doi:10.1038/s41467-019-11526-w}

\bibitem[Kreimer et~al.(2022)]{kreimer2022} Kreimer A, Ashuach T, Inoue F, et al. (2022). Massively parallel reporter perturbation assays uncover temporal regulatory architecture during neural differentiation. \textit{Nature Communications} 13:1504. \href{https://doi.org/10.1038/s41467-022-28659-0}{doi:10.1038/s41467-022-28659-0}

\bibitem[Lee et~al.(2025)]{lee2025} Lee S, McAfee JC, Lee J, et al. (2025). Massively parallel reporter assay investigates shared genetic variants of eight psychiatric disorders. \textit{Cell} 188:1409–1424.e21. \href{https://doi.org/10.1016/j.cell.2024.12.022}{doi:10.1016/j.cell.2024.12.022}

\bibitem[Majoros et~al.(2020)]{majoros2020} Majoros WH, Kim YS, Barrera A, et al. (2020). Bayesian estimation of genetic regulatory effects in high-throughput reporter assays. \textit{Bioinformatics} 36:331–338. \href{https://doi.org/10.1093/bioinformatics/btz545}{doi:10.1093/bioinformatics/btz545}

\bibitem[Matoba et~al.(2026)]{matoba2026} Matoba N, McAfee JC, Krupa O, et al. (2026). Massively parallel assessment of gene regulatory activity at human cortical-structure-associated variants. \textit{Nature Neuroscience}. \href{https://doi.org/10.1038/s41593-026-02454-2}{doi:10.1038/s41593-026-02454-2}

\bibitem[McAfee et~al.(2023)]{mcafee2023} McAfee JC, Lee S, Lee J, et al. (2023). Systematic investigation of allelic regulatory activity of schizophrenia-associated common variants. \textit{Cell Genomics} 3:100404. \href{https://doi.org/10.1016/j.xgen.2023.100404}{doi:10.1016/j.xgen.2023.100404}

\bibitem[Melnikov et~al.(2012)]{melnikov2012} Melnikov A, Murugan A, Zhang X, et al. (2012). Systematic dissection and optimization of inducible enhancers in human cells using a massively parallel reporter assay. \textit{Nature Biotechnology} 30:271–277. \href{https://doi.org/10.1038/nbt.2137}{doi:10.1038/nbt.2137}

\bibitem[Miller et~al.(2014)]{miller2014} Miller JA, Ding SL, Sunkin SM, et al. (2014). Transcriptional landscape of the prenatal human brain. \textit{Nature} 508:199–206. \href{https://doi.org/10.1038/nature13185}{doi:10.1038/nature13185}

\bibitem[Myint et~al.(2019)]{myint2019} Myint L, Avramopoulos DG, Goff LA, Hansen KD. (2019). Linear models enable powerful differential activity analysis in massively parallel reporter assays. \textit{BMC Genomics} 20:209. \href{https://doi.org/10.1186/s12864-019-5556-x}{doi:10.1186/s12864-019-5556-x}

\bibitem[Ng et~al.(2017)]{ng2017} Ng B, White CC, Klein HU, et al. (2017). An xQTL map integrates the genetic architecture of the human brain's transcriptome and epigenome. \textit{Nature Neuroscience} 20:1418–1426. \href{https://doi.org/10.1038/nn.4632}{doi:10.1038/nn.4632}

\bibitem[Liu et~al.(2025)]{liu2025} Liu et al. (2025). Predicting the DNA binding specificity of transcription factor mutants using family-level biophysically interpretable machine learning. \textit{Nucleic Acids Research} 53:gkaf831. \href{https://doi.org/10.1093/nar/gkaf831}{doi:10.1093/nar/gkaf831}

\bibitem[Nielsen et~al.(2016)]{nielsen2016} Nielsen AAK, Der BS, Shin J, et al. (2016). Genetic circuit design automation. \textit{Science} 352:aac7341. \href{https://doi.org/10.1126/science.aac7341}{doi:10.1126/science.aac7341}

\bibitem[NIST(2015)]{nist2015} NIST (2015). \textit{Secure Hash Standard (SHS)}, FIPS 180-4. \url{https://csrc.nist.gov/pubs/fips/180-4/upd1/final}

\bibitem[Peterson et~al.(2012)]{peterson2012} Peterson KA, Nishi Y, Ma W, et al. (2012). Neural-specific Sox2 input and differential Gli-binding affinity provide context and positional information in Shh-directed neural patterning. \textit{Genes \& Development} 26:2802–2816. \href{https://doi.org/10.1101/gad.207142.112}{doi:10.1101/gad.207142.112}

\bibitem[Pollard et~al.(2010)]{pollard2010} Pollard KS, Hubisz MJ, Rosenbloom KR, Siepel A. (2010). Detection of nonneutral substitution rates on mammalian phylogenies. \textit{Genome Research} 20:110–121. \href{https://doi.org/10.1101/gr.097857.109}{doi:10.1101/gr.097857.109}

\bibitem[Ramasamy et~al.(2014)]{ramasamy2014} Ramasamy A, Trabzuni D, Guelfi S, et al. (2014). Genetic variability in the regulation of gene expression in ten regions of the human brain. \textit{Nature Neuroscience} 17:1418–1428. \href{https://doi.org/10.1038/nn.3801}{doi:10.1038/nn.3801}

\bibitem[Rauluseviciute et~al.(2024)]{rauluseviciute2024} Rauluseviciute I, Riudavets-Puig R, Blanc-Mathieu R, et al. (2024). JASPAR 2024: 20th anniversary of the open-access database of transcription factor binding profiles. \textit{Nucleic Acids Research} 52:D174–D182. \href{https://doi.org/10.1093/nar/gkad1059}{doi:10.1093/nar/gkad1059}

\bibitem[Roadmap Epigenomics Consortium(2015)]{roadmap2015} Roadmap Epigenomics Consortium, Kundaje A, Meuleman W, et al. (2015). Integrative analysis of 111 reference human epigenomes. \textit{Nature} 518:317–330. \href{https://doi.org/10.1038/nature14248}{doi:10.1038/nature14248}

\bibitem[Rummel et~al.(2023)]{rummel2023} Rummel CK, Gagliardi M, Ahmad R, et al. (2023). Massively parallel functional dissection of schizophrenia-associated noncoding genetic variants. \textit{Cell} 186:5165–5182.e33. \href{https://doi.org/10.1016/j.cell.2023.09.015}{doi:10.1016/j.cell.2023.09.015}

\bibitem[Rundgren et~al.(2020)]{rundgren2020} Rundgren A, Jordan B, Erdtman S. (2020). \textit{JSON Canonicalization Scheme (JCS)}, RFC 8785. \url{https://www.rfc-editor.org/rfc/rfc8785}

\bibitem[Salomon et~al.(2026)]{salomon2026} Salomon K, Deng C, Dash PM, et al. (2026). Massively parallel characterization and predictive modelling of neuronal regulatory variation. \textit{bioRxiv} (preprint, posted 17 July 2026). \href{https://doi.org/10.64898/2026.07.16.738760}{doi:10.64898/2026.07.16.738760}

\bibitem[Schwartzentruber et~al.(2018)]{schwartzentruber2018} Schwartzentruber J, Foskolou S, Kilpinen H, et al. (2018). Molecular and functional variation in iPSC-derived sensory neurons. \textit{Nature Genetics} 50:54–61. \href{https://doi.org/10.1038/s41588-017-0005-8}{doi:10.1038/s41588-017-0005-8}

\bibitem[Seitz et~al.(2024)]{seitz2024} Seitz EE, McCandlish DM, Kinney JB, Koo PK. (2024). Interpreting cis-regulatory mechanisms from genomic deep neural networks using surrogate models. \textit{Nature Machine Intelligence} 6:701–713. \href{https://doi.org/10.1038/s42256-024-00851-5}{doi:10.1038/s42256-024-00851-5}

\bibitem[Siraj et~al.(2026)]{siraj2026} Siraj L, et al. (2026). Functional dissection of complex trait variants at single-nucleotide resolution. \textit{Nature}. \href{https://doi.org/10.1038/s41586-026-10121-6}{doi:10.1038/s41586-026-10121-6}

\bibitem[Tareen et~al.(2022)]{tareen2022} Tareen A, Kooshkbaghi M, Posfai A, et al. (2022). MAVE-NN: learning genotype-phenotype maps from multiplex assays of variant effect. \textit{Genome Biology} 23:98. \href{https://doi.org/10.1186/s13059-022-02661-7}{doi:10.1186/s13059-022-02661-7}

\bibitem[Tewhey et~al.(2016)]{tewhey2016} Tewhey R, Kotliar D, Park DS, et al. (2016). Direct identification of hundreds of expression-modulating variants using a multiplexed reporter assay. \textit{Cell} 165:1519–1529. \href{https://doi.org/10.1016/j.cell.2016.04.027}{doi:10.1016/j.cell.2016.04.027}

\bibitem[Traets et~al.(2021)]{traets2021} Traets JJ, van der Burght SN, Rademakers S, et al. (2021). Mechanism of life-long maintenance of neuron identity despite molecular fluctuations. \textit{eLife} 10:e66955. \href{https://doi.org/10.7554/eLife.66955}{doi:10.7554/eLife.66955}

\bibitem[Vaishnav et~al.(2022)]{vaishnav2022} Vaishnav ED, de Boer CG, Molinet J, et al. (2022). The evolution, evolvability and engineering of gene regulatory DNA. \textit{Nature} 603:455–463. \href{https://doi.org/10.1038/s41586-022-04506-6}{doi:10.1038/s41586-022-04506-6}

\bibitem[Walker et~al.(2019)]{walker2019} Walker RL, Ramaswami G, Hartl C, et al. (2019). Genetic control of expression and splicing in developing human brain informs disease mechanisms. \textit{Cell} 179:750–771.e22. \href{https://doi.org/10.1016/j.cell.2019.09.021}{doi:10.1016/j.cell.2019.09.021}

\bibitem[Wang et~al.(2020)]{wang2020} Wang G, Sarkar A, Carbonetto P, Stephens M. (2020). A simple new approach to variable selection in regression, with application to genetic fine mapping. \textit{Journal of the Royal Statistical Society Series B} 82:1273–1300. \href{https://doi.org/10.1111/rssb.12388}{doi:10.1111/rssb.12388}

\bibitem[Weiss et~al.(2021)]{weiss2021} Weiss CV, Harshman L, Inoue F, et al. (2021). The cis-regulatory effects of modern human-specific variants. \textit{eLife} 10:e63713. \href{https://doi.org/10.7554/eLife.63713}{doi:10.7554/eLife.63713}

\bibitem[Wilson et~al.(2016)]{wilson2016} Wilson ML, et al. (2016). Genotype Specification Language. \textit{ACS Synthetic Biology} 5:471–478. \href{https://doi.org/10.1021/acssynbio.5b00194}{doi:10.1021/acssynbio.5b00194}

\bibitem[Young et~al.(2021)]{young2021} Young AMH, Kumasaka N, Calvert F, et al. (2021). A map of transcriptional heterogeneity and regulatory variation in human microglia. \textit{Nature Genetics} 53:861–868. \href{https://doi.org/10.1038/s41588-021-00875-2}{doi:10.1038/s41588-021-00875-2}

\bibitem[Zhou and Troyanskaya(2015)]{zhou2015} Zhou J, Troyanskaya OG. (2015). Predicting effects of noncoding variants with deep learning–based sequence model. \textit{Nature Methods} 12:931–934. \href{https://doi.org/10.1038/nmeth.3547}{doi:10.1038/nmeth.3547}

\bibitem[Zubler and Douglas(2009)]{zubler2009} Zubler F, Douglas R. (2009). A framework for modeling the growth and development of neurons and networks. \textit{Frontiers in Computational Neuroscience} 3:25. \href{https://doi.org/10.3389/neuro.10.025.2009}{doi:10.3389/neuro.10.025.2009}

\bibitem[Zubler et~al.(2011)]{zubler2011} Zubler F, Hauri A, Pfister S, et al. (2011). An instruction language for self-construction in the context of neural networks. \textit{Frontiers in Computational Neuroscience} 5:57. \href{https://doi.org/10.3389/fncom.2011.00057}{doi:10.3389/fncom.2011.00057}

\bibitem[Zubler et~al.(2013)]{zubler2013} Zubler F, Hauri A, Pfister S, et al. (2013). Simulating cortical development as a self constructing process: a novel multi-scale approach combining molecular and physical aspects. \textit{PLoS Computational Biology} 9:e1003173. \href{https://doi.org/10.1371/journal.pcbi.1003173}{doi:10.1371/journal.pcbi.1003173}

\end{thebibliography}


\appendix

\section{Процедура с предварително фиксиран план}\label{app:protocol}

Всеки тест в приложение~\ref{app:record} следва една и съща процедура.
\begin{enumerate}
\item \textbf{План.} Файл с план фиксира въпроса, съвкупностите, моделите, базовите модели, мерките, бутстрапа, критериите за решение и етикета, който би получил успешният тест.
\item \textbf{Хеш.} \code{FREEZE.json} записва SHA-256 хешове на плана, скриптовете и входните данни, преди да бъде прочетен какъвто и да е резултат. Прогнозите идват от скрипт, който никога не вижда отделените резултати, и хешът им се добавя към записа.
\item \textbf{Проверка на оценителя.} Скриптът за оценка се изпълнява върху синтетичен положителен случай, при който етикетът трябва да излезе верен, и върху синтетичен нулев случай, при който не бива. Оценител, който не издържи някой от двата, се поправя преди оценката.
\item \textbf{Еднократна оценка.} Оценителят проверява хешовете, оценява запечатаните прогнози и записва един резултат. Променен модел изисква нов план и, ако данните са изразходвани, нови данни.
\item \textbf{Преценка.} Етикетът се присъжда само ако всички критерии са изпълнени. Например долната граница на кандидата е над нулата, кандидатът надминава най-добрия фиксиран базов модел и надминава модел, обучен върху други данни.
\end{enumerate}

Единица за преизвадка е семейството от последователности, генът или достоверният набор (credible set). Единиците се вземат повторно като цели копия в 2000 бутстрап извадки. Интервалите са от 5-ия до 95-ия персентил, така че долният им край е едностранна 95\% граница. Корелациите са корелации на Пиърсън, претеглени по семейства, като всяко изследване има равно тегло. Качеството на подбора е знакова полезност. Това е наблюдаваният ефект, със знака на предсказаната посока, на промените, които моделът избира (по една на семейство), за панел от 10, 50 или 200. При случаен подбор стойността е около нула. Посоката се оценява като точност, като равенствата се броят за половин.

Данните са в четири класа. \emph{Запечатаните} не са отваряни (дялът RESERVE на Siraj). \emph{Изразходваните} са оценени веднъж и не могат да потвърдят по-късни модели (потвържденията на Salomon и на Kosicki). \emph{Разкритите} са насочвали разработката (дялът TUNING на Siraj, по-рано оценените репортерни изследвания и всеки тестов набор с eQTL след първото му оценяване). Резултатите върху тях носят резерва за адаптивност. \emph{Новите} са непознати за всички модели в теста.

\section{Изследователски запис: тестове за пренос от последователността}\label{app:record}

Тестовете питат дали само последователността предсказва какво прави невиждана промяна в ДНК в нервни клетки. Те са доказателството зад границата на преноса от раздел~\ref{sec:transfer}.

\subsection{Задача, данни и модели}

\textbf{Задача.} При дадена референтна и променена последователност да се предскаже репортерният ефект, алтернативен минус референтен (логаритъм с основа~2 от отношение). За естествени варианти да се предскаже посоката на ефекта върху генната експресия.

\textbf{Данни за обучение.} Основният набор е репортерният скрининг в SK-N-SH на \citet{siraj2026} (раздел~\ref{sec:data}). Данните от неврони WTC11-NGN2 \citep{salomon2026} дават активност и ефекти за един клас модели. Във всяко изследване резултатите се центрират и мащабират до единично стандартно отклонение, претеглено по семейства, преди обединяването.

\textbf{Данни за оценка.} Шест нервни репортерни изследвания с алелни ефекти: Lee et al. (GSE244011; 5916 варианта; \citealt{lee2025}), McAfee et al. (GSE211045; 1682; \citealt{mcafee2023}), Matoba et al. (4013; \citealt{matoba2026}), Rummel et al. (GSE225815; 611; \citealt{rummel2023}), Guo et al. (GSE221384; 748; \citealt{guo2023}) и Weiss et al. (GSE152404; 9247; \citealt{weiss2021}). Наборът на Matoba е обучителният дял на проекта. Раздел~\ref{sec:data} описва отделените набори, а таблица~\ref{tab:transfer} дава размерите им. Семейство е свързан компонент от варианти в рамките на 2000 бази (геномен локус за Matoba).

\textbf{Модели.} \emph{k5} е гребенова (ridge) регресия върху разликата в броя на каноничните 5-мери (512 коефициента). \emph{S1} е k5, обучен върху скрининга на Siraj. \emph{S2} добавя разликата в потенциала на мотивите от JASPAR 2024 \citep{rauluseviciute2024} и служи за еталон. \emph{Huber-linear} напасва същите 512 коефициента за думи и 532 за мотиви със загуба на Хюбер и е крайният модел. \emph{S3} и \emph{H3} добавят квадратичен член за контекста на мотивите. Обучени са съответно върху Siraj и върху другите изследвания. \emph{H2} е k5 плюс мотиви, обучен върху другите изследвания. \emph{Cfit} е напасната смес от модели, обучени върху изследванията, и модели, обучени върху Siraj. Други модели са гребенова регресия върху 611 оценки на Sei \citep{chen2022} и конволюционна мрежа, обучена върху обединените ефекти. Има още модел за пренос на активност, научена от измервания на WTC11-NGN2 \citep[вдъхновен от][]{seitz2024,tareen2022}. Последният е модел без обучение, който умножава промяната в заетостта на мотива по експресията на транскрипционния фактор в неокортекса между 8-ата и 13-ата седмица след зачатието \citep{kang2011,miller2014}.

\subsection{По-ранни отрицателни експерименти}

Два експеримента извън серията на Siraj завършиха с отрицателен резултат. Първият е предсказвател за невралната индукция на H1. Фиксираният му модел намалява средноквадратичната грешка с 20,97\% спрямо нулев ефект, но не издържа изискваното сравнение с динуклеотиден модел. Вторият е проверка на модела FamilyCode за свързване на хомеодомени \citep{liu2025} върху \emph{Nkx2-2}, която също не издържа.

\subsection{Проверки на измерванията}

\textbf{Прогноза за пропуснати повторения.} В пилотно проучване на нервен анализ пропускахме по едно РНК повторение и прогнозирахме алелния му дял от останалите, с модела BIRD \citep{majoros2020} или чрез обединяване с аналитично бета-биномно разпределение, и оценихме прогнозите с непрекъснатия ранжиран вероятностен резултат (CRPS; \citealt{gneiting2007}). При дванайсет пропускания в четири компонента от обучителния набор средният CRPS е 0,0549 за BIRD и 0,0622 за обединяването, тоест с 11,7\% по-нисък. Условният 99\% симулационен интервал за разликата е от 0,0044 до 0,0098 и обхваща само грешката на Монте Карло за тази подбрана група. Средните на всичките четири компонента се подобряват, а единственото пропускане в полза на обединяването е запазено. BIRD и обединяването се различават и по приори, така че сравнението не отделя вариацията между повторенията като причина. Последващо сравнение с нормализация по ДНК понижава средния CRPS с 15,8\% спрямо обединяването, но апостериорна прогноза за нулев ефект постига по-добър резултат.

\textbf{Прогноза за по-късни отговори.} В анализа на невралната индукция на \citet{kreimer2022} гребеновата регресия върху четири ранни измервания намалява грешката с 24\% спрямо пренасянето на 12-часовата стойност, при 351 задържани семейства, и се изравнява с по-прост контрол чрез осредняване. Два предварително обучени модела за прогноза без допълнително обучение не се справят по-добре. Тест с необработени броеве, който отделя семейства и инфектирани култури, дава средноквадратична грешка 0,3213 срещу 0,3221 за липса на промяна (с 0,26\% по-добре). И четирите интервала за разликата включват нулата.

\textbf{Шум от броенето.} В GSE225815 при 83\% от вариантите единият алел има по-малко от десет РНК прочитания, затова логаритъмът на отношението се определя главно от псевдоброя и от дисбаланса в ДНК броевете. Високото съгласие между половините на повторенията (0,98) идва от общия ДНК брой и нищо не казва за възпроизводима биология: при 124-те добре измерени варианта изходният слой (head), обучен върху Siraj, дава $r=-0{,}09$. Другите три изследвания с броеве имат пренебрежим шум от броенето (медианна поасонова дисперсия под 0,002 $\log_2^2$). Затова махнахме GSE225815 от основната група в по-късните тестове и го запазихме като анализ на чувствителността.

\subsection{Репортерни ефекти}

\textbf{В рамките на един анализ.} Раздел~\ref{sec:internal} дава вътрешното сравнение на Huber-linear с k5. Две по-ранни потвърждения върху изразходвани данни, за по-ранни напасвани предсказватели от тази серия, не успяха. Потвърждението на Kosicki (3840 промени в 50 семейства) даде грешка 0,016958 за фиксирания предсказвател срещу 0,016919 за нулев ефект. Полезността му при десет промени беше −0,038 срещу +0,023 при случаен подбор. Потвърждението на Salomon (7372 промени в 2938 семейства) даде грешка 0,028957 срещу 0,028998 за k5 и полезност при десет промени 0,170 срещу 0,153. Всичките седем точкови сравнения бяха в полза на модела, но всеки интервал включваше нулата.

\textbf{Между изследвания.} Таблица~\ref{tab:cross} изброява тестовете, които прилагат фиксиран модел върху други изследвания. Еталонът S2 е обучен само върху Siraj. Той има $r\approx0{,}047$ (S1 има 0,047 върху всичките шест изследвания, интервал от 0,026 до 0,069), тоест $r^2\approx0{,}002$. S2 е положителен в четири от петте изследвания без GSE225815. След проверката на шума от броенето тези пет изследвания са основната група. Групата от шест се пази като анализ на чувствителността.

\begin{table}[tb]
\small
\caption{Тестове между изследвания. $r$ е корелация на Пиърсън, претеглена по семейства, с равно тегло на изследванията (от 5-ия до 95-ия персентил). Всеки тест е оценен веднъж, на 19 и 20 октомври 2026 г. Нито един тест не получи етикета си.}\label{tab:cross}
\begin{tabular}{@{}r L{62mm} L{34mm} L{28mm}@{}}
\toprule
№ & Кандидат и съвкупност & $r$ на кандидата & $r$ на еталона \\
\midrule
1 & H2, обучен върху другите две изследвания; дялове McAfee и Matoba & 0,015 (−0,024; 0,054) & S1 0,060 \\
2 & Смес Cfit; три новооценени изследвания & 0,046 (0,009; 0,082) & S2 0,050 \\
3 & Гребенова регресия върху Sei; шест изследвания (смес със S2: 0,041) & −0,003 (−0,024; 0,019) & S2 0,047 \\
4 & Квадратичен контекстен слой S3; шест изследвания & 0,046 & S2 0,047 \\
5 & Пренос на активност от WTC11-NGN2; шест изследвания & −0,010 (−0,029; 0,011) & S2 0,047 \\
6 & Обединена конволюционна мрежа; пет изследвания & 0,014 (−0,008; 0,035) & S2 0,047 \\
7 & Свързване на експресирани транскрипционни фактори без обучение; пет изследвания & 0,0015 & S2 0,0465 \\
8 & Huber-linear, фиксиран; пет изследвания & 0,0490 & S2 0,0465 \\
\bottomrule
\end{tabular}
\end{table}

Тест~8 е крайният модел (раздел~\ref{sec:transfer}, който дава и сравнението с прогноза за нулев ефект). Корелациите на S2 по изследвания са от −0,014 (Guo et al.) до 0,115 (Weiss et al.). В проучвателен анализ с подбор по резултата знакът на предсказанието на S2 е верен за 53\% до 57\% от 10-те процента най-големи ефекти в изследванията на Lee, Matoba и Weiss, срещу 50\% общо. Това е само посока за проверка. Еволюционната консервация (phyloP100way; \citealt{pollard2010}) не помага. Корелацията ѝ с абсолютния ефект е средно 0,017, а умножаването на S2 по консервацията понижава средното $r$ от 0,047 на 0,034.

\subsection{Естествени варианти}

\textbf{Посока на ефектите върху експресията.} Данните за обучение бяха QTL с фино картиране от мозък на възрастни: мозъчни региони на GTEx, BrainSeq, CommonMind, ROSMAP и Braineac2 в eQTL Catalogue \citep{gtex2020,jaffe2018,fromer2016,ng2017,ramasamy2014,kerimov2021}. Те съдържат 14\,617 двойки вариант-ген с апостериорна вероятност за включване (PIP) поне 0,5, с фино картиране чрез SuSiE \citep{wang2020}. Валидирането задържа цели хромозоми. Тестовете използваха набори от фетален мозък и от клетки, получени от нервни. M4 използва оценки на хроматина от Sei по интервал на разстоянието до началото на транскрипцията (TSS). M5 добавя промяната в GC, локалния k5 и изходния слой на репортерния модел. Базови модели са B1 (алтернативният алел печели GC), B2 (логистична регресия върху k5) и B3 (изходният слой, обучен върху Siraj, без допълнително напасване).

\begin{table}[tb]
\small
\caption{Точност на посоката върху задържани eQTL с фино картиране; преизвадка по гени, от 5-ия до 95-ия персентил.}\label{tab:direction}
\begin{tabular}{@{}L{54mm}r r r L{42mm}@{}}
\toprule
Тестов набор & $n$ & M4 & M5 & M5 спрямо най-добрия базов \\
\midrule
Фетален неокортекс \citep{walker2019} & 792 & 0,547 & 0,566 & B3: +0,028 (−0,009; +0,065) \\
Прогенитори \citep{aygun2021} & 229 & 0,528 & 0,646 & B2: +0,066 (0,013; 0,117) \\
Неврони \citep{aygun2021} & 80 & 0,575 & 0,600 & широки интервали \\
Диференциращи се неврони \citep{jerber2021} & 320 & 0,534 & 0,603 & B2: +0,003 (−0,037; +0,045) \\
Сетивни неврони \citep{schwartzentruber2018} & 519 & 0,538 & 0,518 & B2 0,499 (PIP ≥ 0,1) \\
Микроглия \citep{young2021} & 272 & няма & 0,570 & изходен слой 0,574 \\
\bottomrule
\end{tabular}
\end{table}

Обединено за четирите набора от клетки с нервен произход при PIP ≥ 0,5 (626 двойки) M5 достига 0,599 (0,566; 0,630) срещу 0,567 за B2. Най-голямата печалба в един набор, при прогениторите, не се повтори при наборите на Jerber и на сетивните неврони. Нито един модел не стои последователно над базовия модел с k-мери. Трите етикета (за посоката във фетален мозък, за потвърждение в диференциращи се неврони и за повторение в клетки с нервен произход) не бяха присъдени.

\textbf{Кой вариант в достоверния набор е причинният?} Подреждане на членовете на достоверните набори с условна логит регресия беше обучено върху 5836 набора от мозък на възрастни (27\,875 члена) със задържане на цели хромозоми. Оценихме точността на първото място (top-1) върху 439 набора от фетален неокортекс \citep{walker2019} и 420 набора от клетки с нервен произход (таблица~\ref{tab:ranking}).

\begin{table}[tb]
\small
\caption{Точност на първото място при подреждане в достоверен набор (от 5-ия до 95-ия персентил).}\label{tab:ranking}
\begin{tabular}{@{}L{62mm}L{40mm}L{40mm}@{}}
\toprule
Метод & Фетален неокортекс & Набори от клетки с нервен произход \\
\midrule
Случаен член & 0,355 & 0,328 \\
Най-близък TSS (B1) & \textbf{0,460} (0,420; 0,499) & 0,402 (0,362; 0,442) \\
Интервали на разстоянието (B2) & 0,433 (0,403; 0,463) & 0,381 (0,351; 0,411) \\
Най-висок phyloP (B3) & 0,355 (0,320; 0,392) & 0,333 (0,294; 0,373) \\
Разстояние + Sei (M4) & 0,435 (0,396; 0,474) & \textbf{0,426} (0,386; 0,464) \\
Разстояние + Sei + phyloP (M5) & 0,419 (0,381; 0,458) & 0,421 (0,381; 0,461) \\
\bottomrule
\end{tabular}
\end{table}

Никой модел не надминава правилото за най-близкия TSS във фетален неокортекс. При наборите от клетки с нервен произход M4 надминава интервалите на разстоянието с +0,045 (+0,005; +0,085). Той не надминава ясно B1 (M5 − B1 = +0,019, интервал от −0,029 до +0,067). Консервацията не носи сигнал: phyloP избира члена с най-висока вероятност от фино картиране не по-добре от случайността (AUC по двойки 0,495 в обучителните набори и 0,503 при Walker). Добавянето на достъпност на хроматина по DNase от Roadmap \citep{roadmap2015} не помага във фетален неокортекс. Двата модела с достъпност достигат 0,417 и 0,420, под M4 (0,435) и под правилото за най-близкия TSS (0,460). Нито един от двата етикета не беше присъден.

\subsection{Обобщение}

В тринайсетте теста изходният слой, обучен върху Siraj, носи слаб, устойчив между изследванията сигнал ($r$ около 0,05), към който нищо от опитаното не добавя. Сигнали, които изглеждат по-добре в един набор, не се повтарят в следващия. Простите базови модели (най-близък TSS, броят на k-мери, изходният слой на репортерния модел) са трудни за надминаване при естествени варианти.

\section{Фиксирани тестове и къде се намират}\label{app:tests}

Всяка папка в \code{research/frozen-tests/} съдържа плана (\code{PLAN.md}), хешовете (\code{FREEZE.json}), \code{scripts/}, малки файлове с резултати и \code{RESULTS.md}. Оригиналното дърво е \code{build/neural-predictor-reviewed-20261005/} и не е под версионен контрол.

\begin{center}\small
\begin{tabular}{@{}r L{70mm} L{62mm}@{}}
\toprule
№ & Тест & Папка \\
\midrule
1 & Изключване на по едно изследване, три изследвания & \code{hnp-loso-20261019} \\
2 & Комбиниран предсказвател, шест изследвания & \code{hnp-loso6-20261019} \\
3 & Гребенова регресия върху Sei & \code{sei-features-20261019} \\
4 & Квадратичен контекстен слой & \code{context-loso6-20261019} \\
5 & Модел за пренос на активност & \code{activity-transfer-loso6-20261019} \\
6 & Обединена конволюционна мрежа & \code{cnn-loso6-20261019} \\
7 & Свързване на експресирани транскрипционни фактори без обучение & \code{mechanism-zeroshot-20261019} \\
8 & Пренос на Huber-linear & \code{huber-transfer-20261020} \\
9 & Посока на eQTL, фетален неокортекс & \code{eqtl-brain-20261019} \\
10 & Посока на eQTL, Jerber & \code{eqtl-brain-20261019} (\code{PLAN\_CONFIRM.md}) \\
11 & Посока на eQTL, сетивни неврони и микроглия & \code{eqtl-brain-20261019} (\code{PLAN\_REP.md}) \\
12 & Подреждане в достоверен набор & \code{eqtl-causal-20261019} (\code{PLAN.md}) \\
13 & Подреждане в достоверен набор с достъпност & \code{eqtl-causal-20261019} (\code{PLAN2.md}) \\
\bottomrule
\end{tabular}
\end{center}

Проучвателен, нефиксиран анализ на консервацията и посоката е в \code{constraint-factor-20261019}.

\end{document}
